% ===== PRÉAMBULE CANONIQUE — à coller VERBATIM en tête de CHAQUE .tex =====
\documentclass[11pt,a4paper]{article}
\usepackage[utf8]{inputenc}\usepackage[T1]{fontenc}\usepackage{lmodern}
\usepackage[french]{babel}
\usepackage{amsmath,amssymb,mathtools}\usepackage{icomma}
\usepackage{siunitx}\sisetup{locale=FR}  % \qty{}{}, \si{}, \ang{}
\usepackage{xcolor}\usepackage{colortbl}
\usepackage{tikz}\usepackage{pgfplots}\pgfplotsset{compat=1.18}
\usepackage[siunitx,europeanresistors]{circuitikz} % circuits (élec.)
\usepackage[most]{tcolorbox}\usepackage{enumitem}\usepackage{array,booktabs}
\usepackage[a4paper,margin=2.2cm]{geometry}
\usetikzlibrary{arrows.meta,calc,angles,quotes,positioning,patterns,%
  backgrounds,shapes.geometric,decorations.pathmorphing,decorations.markings,intersections}
\definecolor{primary}{RGB}{222,49,99}\definecolor{primarydark}{RGB}{150,22,55}
\definecolor{defcol}{RGB}{0,114,178}\definecolor{methcol}{RGB}{52,92,148}
\definecolor{excol}{RGB}{0,135,90}\definecolor{attncol}{RGB}{200,40,55}
\tcbset{boxbase/.style={enhanced,breakable,boxrule=0.9pt,arc=2.5pt,
  left=3mm,right=3mm,top=2.3mm,bottom=2.3mm,fonttitle=\bfseries,coltitle=white}}
% --- boîtes TUTORIEL (noms EXACTS, ne pas renommer) ---
\newtcolorbox{cle}[1][]{boxbase,colback=defcol!6,colframe=defcol,
  title={\textsc{À retenir}\ifx\relax#1\relax\relax\else\textmd{ : \itshape #1}\fi}}
\newtcolorbox{methode}[1][]{boxbase,colback=methcol!7,colframe=methcol,sharp corners,
  title={\textsc{Méthode}\ifx\relax#1\relax\relax\else\textmd{ --- \itshape #1}\fi}}
\newtcolorbox{exemplebox}[1][]{boxbase,colback=excol!5,colframe=excol,
  title={\textsc{Exemple}\ifx\relax#1\relax\relax\else\textmd{ : \itshape #1}\fi}}
\newtcolorbox{piege}[1][]{boxbase,colback=attncol!6,colframe=attncol,
  title={\textsc{Piège}\ifx\relax#1\relax\relax\else\textmd{ : \itshape #1}\fi}}
\newtcolorbox{remarque}[1][]{boxbase,colback=black!4,colframe=black!45,
  title={\textsc{Remarque}\ifx\relax#1\relax\relax\else\textmd{ : \itshape #1}\fi}}
% --- boîtes EXERCICE / CORRIGÉ (noms EXACTS, auto-comptées) ---
\newtcolorbox[auto counter]{exo}[1][]{boxbase,colback=primary!4,colframe=primary,
  title={\textsc{Exercice}~\thetcbcounter\ifx\relax#1\relax\relax\else\textmd{ --- \itshape #1}\fi}}
\newtcolorbox[auto counter]{corrige}[1][]{boxbase,colback=excol!4,colframe=excol,
  title={\textsc{Corrigé}~\thetcbcounter\ifx\relax#1\relax\relax\else\textmd{ --- \itshape #1}\fi}}
\newcommand{\vv}[1]{\vec{#1}}\newcommand{\ds}{\displaystyle}
\newcommand{\R}{\mathbb{R}}\newcommand{\N}{\mathbb{N}}\newcommand{\Z}{\mathbb{Z}}
% (un \usepackage{fancyhdr}+en-tête est possible ; il est ignoré côté web)
% =========================================================================
\begin{document}

\begin{center}
{\large\bfseries\color{excol} Thème 1 — Corrigé Difficile}\\[-2pt]
{\color{excol}\rule{5cm}{1pt}}
\end{center}

\begin{corrige}[la longueur d'onde à partir d'un enregistrement]
Le graphe est en fonction du \emph{temps} : on y lit $T$. La longueur d'onde vient ensuite de
$\lambda=vT$.
\begin{enumerate}[label=\alph*)]
  \item Le motif se répète tous les \SI{2}{\second} : $T=\SI{2}{\second}$, donc
        $f=\dfrac{1}{T}=\SI{0.5}{\hertz}.$
  \item $\lambda=vT=\SI{0.5}{\metre\per\second}\times\SI{2}{\second}=\SI{1}{\metre}=\SI{100}{\centi\metre}.$
  \item \textbf{Non.} Un graphe $y(t)$ ne montre que le mouvement dans le temps d'\emph{un
        seul} point : on y lit $T$, pas $\lambda$. Pour lire $\lambda$ directement, il faudrait
        une \emph{photo} de la corde, c.-à-d. un graphe $y(x)$.
\end{enumerate}
\end{corrige}

\begin{corrige}[deux graphes, une vitesse]
Chaque graphe donne une seule information ; on les combine avec $v=\lambda/T$.
\begin{enumerate}[label=\alph*)]
  \item La longueur d'onde se lit sur la \emph{photo} $y(x)$ (axe en mètres) : le motif se
        répète tous les \SI{4}{\metre}, donc $\lambda=\SI{4}{\metre}.$
  \item La période se lit sur le graphe du \emph{point fixe} $y(t)$ (axe en secondes) : le
        motif se répète tous les \SI{0.5}{\second}, donc $T=\SI{0.5}{\second}.$
  \item $v=\dfrac{\lambda}{T}=\dfrac{\SI{4}{\metre}}{\SI{0.5}{\second}}=\SI{8}{\metre\per\second}.$
\end{enumerate}
\end{corrige}

\begin{corrige}[échographie : la profondeur d'un obstacle]
On lit la durée sur l'oscilloscope, puis on n'oublie pas que l'onde fait l'\emph{aller-retour}.
\begin{enumerate}[label=\alph*)]
  \item 8 divisions à \SI{10}{\micro\second} chacune :
        $\Delta t = 8\times\SI{10}{\micro\second}=\SI{80}{\micro\second}=\SI{80e-6}{\second}.$
  \item Pendant $\Delta t$, l'onde parcourt l'aller \emph{et} le retour, soit $2d$ :
        \[ 2d = v\,\Delta t \ \Rightarrow\ d=\frac{v\,\Delta t}{2}
           =\frac{\SI{1500}{\metre\per\second}\times\SI{80e-6}{\second}}{2}
           =\frac{\SI{0.12}{\metre}}{2}=\SI{0.06}{\metre}=\SI{6}{\centi\metre}. \]
        \textbf{Piège :} oublier le facteur $2$ donnerait \SI{12}{\centi\metre}, soit le double.
\end{enumerate}
\end{corrige}

\begin{corrige}[rides à la surface de l'eau]
Deux crêtes successives sont distantes d'une longueur d'onde $\lambda$ : on compte les crêtes
en partant de $S$.
\begin{enumerate}[label=\alph*)]
  \item La 1\textsuperscript{re} crête est à $\lambda$ de $S$, la 2\textsuperscript{e} à $2\lambda$,
        la 3\textsuperscript{e} à $3\lambda$. Donc $d = SA = 3\lambda.$
  \item $d = 3\lambda = 3\times\SI{4}{\centi\metre}=\SI{12}{\centi\metre}.$
\end{enumerate}
\end{corrige}

\begin{corrige}[un son qui change de milieu]
La fréquence est fixée par la source et reste invariante ; $v$ et $\lambda$ changent avec le milieu.
\begin{enumerate}[label=\alph*)]
  \item $\lambda_{\text{eau}}=\dfrac{v_{\text{eau}}}{f}
        =\dfrac{\SI{1500}{\metre\per\second}}{\SI{250}{\hertz}}=\SI{6}{\metre}.$
  \item La fréquence \textbf{ne change pas} : elle est imposée par la source. On garde
        $f=\SI{250}{\hertz}$ dans l'air.
  \item $\lambda_{\text{air}}=\dfrac{v_{\text{air}}}{f}
        =\dfrac{\SI{340}{\metre\per\second}}{\SI{250}{\hertz}}=\SI{1.36}{\metre}.$
        Rapport : $\dfrac{\lambda_{\text{eau}}}{\lambda_{\text{air}}}=\dfrac{6}{1{,}36}\approx 4{,}4$,
        exactement le rapport des célérités $\dfrac{1500}{340}\approx 4{,}4$.
\end{enumerate}
\end{corrige}

\end{document}
